\chapter[PENDAHULUAN]{Pendahuluan}
\label{chap:pendahuluan}

\section{Latar Belakang}


Permasalahan dalam pengelolaan pengetahuan tidak selalu terletak pada ketiadaan
informasi. Sebuah pertanyaan pengguna sering kali memiliki jawaban yang telah
tersedia dalam dokumen rujukan. Namun, jumlah, panjang, dan struktur dokumen yang
beragam dapat menyebabkan informasi tersebut tersebar di berbagai bagian atau
bahkan di beberapa dokumen yang saling berkaitan. Oleh karena itu, pengguna tidak
hanya perlu menemukan dokumen yang relevan, tetapi juga mengidentifikasi dan
memahami konteksnya, serta menghubungkan informasi yang tersebar untuk memperoleh
jawaban yang utuh dan tepat.

Skala dan kompleksitas dokumen membuat pencarian informasi menjadi tantangan pada
berbagai konteks. Dalam hukum, JDIHN mencatat ratusan ribu peraturan
perundang-undangan yang memiliki struktur dan hubungan berbeda, termasuk hubungan
hierarki, rujukan, perubahan, dan pencabutan
\parencite{bphnJDIHNDokumenHukum2025,uuPembentukanPeraturan2011}. Kesalahan dalam
menemukan atau memahami hubungan tersebut dapat menyebabkan penggunaan dasar
hukum yang tidak tepat \parencite{faisalGranularityAwareLegal2024}. Tantangan
serupa muncul pada \emph{customer service} ketika pertanyaan pengguna membutuhkan
informasi produk dan prosedur yang tersebar di berbagai dokumen. Petugas \emph{customer service}
perlu menelusuri kembali dokumen tersebut untuk memberikan jawaban yang
sesuai \parencite{liQuestionAnsweringTechnical2018}.

Perkembangan LLM dapat membantu pengguna memperoleh informasi tanpa harus memahami struktur
atau kata kunci yang digunakan dalam dokumen. Meskipun demikian, LLM masih
memiliki keterbatasan karena jawaban yang dihasilkan bergantung pada pengetahuan
parametris yang tersimpan di dalam model. Informasi tersebut tidak selalu
terbaru dan dapat diverifikasi sehingga dapat menyebabkan
\emph{hallucination} \parencite{jiSurveyHallucinationNatural2023}. Pendekatan RAG kemudian
digunakan untuk mengatasi keterbatasan tersebut dengan mengambil informasi yang
relevan dari sumber eksternal sebelum jawaban dibentuk oleh LLM
\parencite{lewisRetrievalAugmentedGenerationKnowledgeIntensive2021}.

Namun, LLM tidak selalu mampu memanfaatkan seluruh informasi yang tersedia dalam
dokumen eksternal. Kapasitas \emph{context window} membatasi jumlah informasi yang
dapat diproses dalam satu permintaan. Selain itu, penambahan konteks yang lebih
panjang tidak selalu meningkatkan kualitas jawaban karena model dapat mengabaikan
informasi penting yang berada di dalam konteks tersebut
\parencite{liuLostMiddle2024}. Oleh karena itu, \emph{pipeline} RAG diperlukan untuk
mengolah dokumen, menemukan bagian yang relevan, dan memberikan bagian tersebut
kepada LLM sebagai konteks
\parencite{lewisRetrievalAugmentedGenerationKnowledgeIntensive2021}.

Pemrosesan utama tersebut diawali dengan \emph{chunking}, yaitu membagi dokumen
menjadi beberapa potongan informasi. Penentuan ukuran dan batas \emph{chunk}
menjadi tantangan karena \emph{chunk} yang terlalu besar dapat memuat informasi
yang tidak relevan dan menimbulkan \emph{noise}, sedangkan \emph{chunk} yang
terlalu kecil dapat kehilangan konteks dan memutus makna semantis informasi.
Pemotongan pada batas yang tidak tepat juga dapat memisahkan definisi,
pengecualian, atau rujukan yang seharusnya dipahami secara bersamaan. Akibatnya,
informasi yang benar dapat lebih sulit ditemukan atau tidak cukup lengkap untuk
digunakan oleh LLM. Penelitian NitiBench menunjukkan bahwa teknik yang
mempertimbangkan hierarki dan rujukan silang pada dokumen hukum dapat memberikan
hasil yang berbeda dibandingkan pendekatan yang tidak memperhatikan struktur
dokumen \parencite{akarajaradwongNitiBenchComprehensiveStudy2025}. Oleh karena itu,
\emph{chunker} perlu mempertimbangkan struktur dan karakteristik dokumen, bukan
hanya ukuran karakter atau jumlah \emph{token}.

Setelah \emph{chunk} terbentuk, sistem membutuhkan \emph{indexer} agar potongan
informasi yang relevan dapat ditemukan secara efisien. Pencarian berbasis kata
dapat membantu menemukan istilah yang sama persis dengan pertanyaan, sedangkan
pencarian berbasis vektor dapat menemukan informasi yang memiliki kemiripan makna
meskipun menggunakan kata yang berbeda
\parencite{robertsonProbabilisticRelevanceFramework2009,
karpukhinDensePassageRetrieval2020}. Kedua pendekatan tersebut memiliki
keterbatasan. Pencarian berbasis kata dapat melewatkan informasi yang
menggunakan istilah berbeda, sementara pencarian berbasis vektor dapat mengambil
potongan yang memiliki makna umum tetapi tidak cukup spesifik. Selain itu,
penelitian mengenai \emph{retrieval granularity} menunjukkan bahwa pilihan unit
yang diindeks, seperti dokumen, paragraf, kalimat, atau proposisi, dapat
memengaruhi performa pencarian dan kualitas jawaban yang dihasilkan
\parencite{chenDenseXRetrieval2023}. Performa sistem \emph{retrieval} juga dapat berubah
ketika model diterapkan pada \emph{domain} yang berbeda. BEIR menunjukkan bahwa model
\emph{retrieval} perlu diuji pada berbagai jenis tugas dan koleksi dokumen karena
performanya pada satu \emph{domain} tidak selalu mencerminkan performanya pada \emph{domain}
lain \parencite{thakurBEIRHeterogeneousBenchmark2021}. Dengan demikian,
\emph{indexer} perlu dirancang dengan mempertimbangkan strategi pencarian, unit
informasi, model representasi, dan karakteristik \emph{domain}.

Meskipun \emph{indexer} dapat menemukan \emph{chunk} yang relevan, pencarian
berbasis teks belum selalu mampu menghubungkan informasi yang tersebar di
berbagai bagian atau dokumen. \emph{Knowledge graph} digunakan untuk
merepresentasikan entitas dan hubungan antarinformasi secara eksplisit sehingga
sistem dapat menelusuri keterkaitan, menggabungkan informasi dari beberapa
\emph{chunk}, dan mendukung pertanyaan yang membutuhkan penalaran bertahap. Pada
dokumen hukum, graf dapat merepresentasikan hierarki, rujukan silang, perubahan,
dan pencabutan peraturan. Pada \emph{domain} lain, graf dapat menghubungkan produk,
fitur, prosedur, dan kategori \emph{service}. Namun, pembangunan graf memiliki
tantangan karena struktur hubungan berbeda antar-\emph{domain}, proses ekstraksi
dapat menghasilkan relasi yang keliru, dan perubahan dokumen perlu diperbarui
tanpa menghilangkan keterlacakan terhadap sumber asalnya
\parencite{edgeFromLocalGlobalGraph2024,
xuRetrievalAugmentedGenerationKnowledgeGraphs2024,
zhongComprehensiveSurveyAutomatic2023}.

Permasalahan tersebut tidak hanya muncul ketika sistem pertama kali dibangun,
tetapi juga berulang setiap kali sumber pengetahuan atau \emph{domain} baru
ditambahkan. Setiap \emph{domain} dapat memiliki struktur, format, istilah, dan
hubungan informasi yang berbeda sehingga strategi \emph{chunking}, pengindeksan,
serta konstruksi \emph{knowledge graph} yang sesuai pada satu \emph{domain}
belum tentu memberikan hasil yang sama pada \emph{domain} lain. Perbedaan ini
terlihat pada penelitian \textcite{stablerImpactChunkingDomainSpecific2025}, yang
menunjukkan bahwa konfigurasi \emph{sentence splitting} dengan jendela 512 \emph{token}
dan \emph{overlap} 200 \emph{token} menghasilkan nilai \emph{Intersection-over-Union}
sekitar 0,06--0,11 pada dokumen \emph{finance}, tetapi hanya sekitar 0,03--0,04
pada dokumen \emph{computer science} dan \emph{astrophysics}. Studi lain juga
menunjukkan bahwa \emph{dynamic token chunking} lebih sesuai untuk dokumen
kesehatan dan fisika, sedangkan \emph{paragraph grouping} memberikan hasil yang
lebih baik pada dokumen hukum dan matematika
\parencite{shaukatChunkingEmbedding2026}. Temuan tersebut menunjukkan bahwa
konfigurasi \emph{chunking} tidak dapat ditentukan secara universal. Jika seluruh
komponen terikat langsung pada kode inti, penambahan \emph{domain} baru akan
menuntut perubahan kode, pemrosesan ulang dokumen, pembangunan ulang \emph{artifact}, dan
pengujian ulang. Oleh karena itu, diperlukan \emph{sistem plugin} yang
memungkinkan \emph{chunker}, \emph{indexer}, dan pembangun \emph{knowledge graph}
dipilih atau ditukar sesuai kebutuhan melalui kontrak masukan dan keluaran yang
jelas. Sistem tersebut juga perlu mempertahankan identitas, versi, kapabilitas,
aturan kompatibilitas, dan keterlacakan \emph{artifact} agar fleksibilitas tetap menjaga
konsistensi proses \parencite{hanRetrievalAugmentedGenerationGraphs2025,
gaoModularRAG2024,garlanIntroductionSoftware1994}.

Pada tahap percakapan, masukan pengguna dan konteks hasil \emph{retrieval} perlu
dijaga. \emph{Guardrail} diperlukan karena keduanya dapat mengandung instruksi
yang tidak seharusnya memengaruhi perilaku model. Penelitian mengenai
\emph{indirect prompt injection} menunjukkan bahwa informasi eksternal dapat
diperlakukan sebagai instruksi oleh LLM dan digunakan untuk mengubah perilaku
aplikasi \parencite{greshakeIndirectPromptInjection2023}. Basis pengetahuan juga
dapat dimanipulasi melalui penyisipan teks berbahaya yang memengaruhi hasil
\emph{retrieval} dan jawaban model \parencite{zouPoisonedRAG2025}. PII dalam kueri,
histori, dan dokumen konteks perlu dilindungi karena RAG dapat membuka risiko
kebocoran data privat \parencite{mccallisterGuideProtectingPII2010,
zengGoodBadPrivacyRAG2024}. LongMemEval menunjukkan bahwa percakapan panjang
dapat menyebabkan sistem kehilangan informasi, sedangkan ringkasan rekursif dapat
mempertahankan informasi dan meningkatkan konsistensi respons
\parencite{wuLongMemEval2025,wangRecursivelySummarizingLongTerm2023}. Karena itu,
sistem memerlukan \emph{guardrail}, perlindungan PII, dan \emph{history compactor}.

\section{Rumusan Masalah}

Berdasarkan latar belakang tersebut, rumusan masalah penelitian adalah sebagai
berikut.

\begin{enumerate}
    \item Bagaimana sistem menerapkan strategi \emph{chunking} untuk membagi dokumen dengan karakteristik beragam menjadi unit informasi, mempertahankan konteks struktural dan hubungan antarbagiannya, serta mendukung penambahan \emph{domain} baru?
    \item Bagaimana sistem menerapkan strategi pengindeksan untuk mengubah unit informasi hasil \emph{chunking} menjadi representasi indeks yang dapat menemukan informasi berdasarkan kesesuaian istilah dan makna terhadap pertanyaan pengguna, serta mendukung penambahan \emph{domain} baru?
    \item Bagaimana sistem menerapkan mekanisme \emph{knowledge graph} untuk merepresentasikan dan menelusuri hubungan antarentitas yang diperlukan dalam menjawab pertanyaan pengguna dari satu atau beberapa dokumen, serta mendukung penambahan \emph{domain} baru?
\end{enumerate}

\section{Tujuan dan Ukuran Keberhasilan Pencapaian}

Tujuan utama penelitian ini adalah mengembangkan komponen pembentukan basis
pengetahuan pada sistem \emph{Retrieval-Augmented Generation} (\emph{RAG})
\emph{multi-domain} yang mencakup \emph{chunker}, \emph{indexer}, dan \emph{knowledge
graph} untuk mengubah dokumen dengan karakteristik beragam menjadi unit
informasi, indeks pencarian, serta representasi hubungan yang dapat digunakan
pada tahap \emph{retrieval}, sekaligus menyediakan arsitektur yang dapat
diperluas ke sumber pengetahuan atau \emph{domain} baru tanpa mengubah kode inti.

Dalam penelitian ini, konstruksi basis pengetahuan didefinisikan sebagai rangkaian
proses pembentukan representasi pengetahuan yang mencakup tiga tahap utama, yaitu
(1) membagi representasi pengetahuan menjadi unit-unit informasi yang lebih kecil
dengan tetap mempertahankan konteksnya, (2) mengindeks atau memproyeksikan unit
tersebut ke dalam \emph{vector database} agar dapat dicari, dan (3) membentuk
entitas, informasi, serta relasi dalam bentuk triplet ke dalam \emph{knowledge graph}.

Secara lebih rinci, tujuan penelitian ini adalah sebagai berikut.

\begin{enumerate}
    \item Mengembangkan strategi \emph{chunking} yang membagi dokumen dengan karakteristik beragam menjadi unit informasi dengan tetap mempertahankan konteks struktural dan hubungan antarbagiannya, serta dapat diperluas ke sumber pengetahuan atau \emph{domain} baru tanpa mengubah kode inti.
    \item Mengembangkan \emph{indexer} yang mengubah unit informasi hasil \emph{chunking} menjadi representasi indeks untuk menemukan informasi berdasarkan kesesuaian istilah dan makna terhadap pertanyaan pengguna, serta dapat digunakan pada sumber pengetahuan atau \emph{domain} baru tanpa mengubah kode inti.
    \item Mengembangkan mekanisme pembentukan \emph{knowledge graph} yang merepresentasikan dan menelusuri hubungan antarentitas dari satu atau beberapa dokumen, serta dapat diperluas ke sumber pengetahuan atau \emph{domain} baru tanpa mengubah kode inti.
\end{enumerate}

Ketercapaian setiap tujuan diukur melalui kriteria berikut.

\begin{enumerate}
    \item Ketercapaian tujuan pertama diukur melalui:
    \begin{enumerate}[label=\alph*.]
        \item Setiap hasil \emph{chunking} memiliki identitas yang jelas dan dapat ditelusuri kembali ke dokumen, bagian dokumen, posisi sumber, serta struktur yang menjadi dasar pembentukannya, termasuk hubungan antar-\emph{chunk} seperti relasi \emph{previous} dan \emph{next}.
        \item \emph{Chunker} menghasilkan unit informasi yang dapat digunakan oleh \emph{indexer} dan \emph{knowledge graph}.
        \item Hasil terbaik dari perbandingan beberapa strategi \emph{chunking} pada \emph{domain} \emph{customer service} Netgear dan dokumen hukum pendidikan Indonesia ditentukan berdasarkan nilai \emph{recall@K}.
        \item Penambahan strategi \emph{chunking} untuk sumber atau \emph{domain} baru dapat dilakukan melalui komponen tambahan tanpa mengubah kode inti.
    \end{enumerate}

    \item Ketercapaian tujuan kedua diukur melalui:
    \begin{enumerate}[label=\alph*.]
        \item Sistem dapat menyediakan beberapa strategi pencarian, yaitu pencarian berbasis istilah atau \emph{sparse retrieval}, pencarian berbasis representasi vektor atau \emph{dense retrieval}, serta kombinasi keduanya melalui \emph{hybrid retrieval}.
        \item Setiap hasil indeks mempertahankan identitas \emph{chunk}, identitas dokumen, \emph{metadata}, konfigurasi, model, dan informasi asal yang diperlukan untuk penelusuran.
        \item Hasil terbaik dari perbandingan strategi pencarian pada \emph{domain} \emph{customer service} Netgear dan dokumen hukum pendidikan Indonesia ditentukan berdasarkan nilai \emph{recall@K}.
        \item Penambahan strategi \emph{indexer} untuk sumber atau \emph{domain} baru dapat dilakukan melalui komponen tambahan tanpa mengubah kode inti.
    \end{enumerate}

    \item Ketercapaian tujuan ketiga diukur melalui:
    \begin{enumerate}[label=\alph*.]
        \item \emph{Knowledge graph} dapat dibangun dari \emph{artifact} \emph{chunk} yang sama dengan yang digunakan oleh \emph{indexer} tanpa melakukan pemotongan ulang atau menggunakan sumber yang berbeda.
        \item \emph{Knowledge graph} menghasilkan hubungan antarentitas yang dapat ditelusuri dari satu atau beberapa dokumen.
        \item Hasil terbaik pembentukan hubungan pada \emph{domain} \emph{customer service} Netgear dan dokumen hukum pendidikan Indonesia ditentukan berdasarkan nilai \emph{recall@K}.
        \item Penambahan mekanisme pembentukan \emph{knowledge graph} untuk sumber atau \emph{domain} baru dapat dilakukan melalui komponen tambahan tanpa mengubah kode inti.
    \end{enumerate}
\end{enumerate}

\section{Batasan Penelitian}

Agar penelitian dapat dilaksanakan secara terarah, ruang lingkup penelitian perlu
dibatasi secara eksplisit. Pembatasan ini bertujuan untuk menjaga fokus penelitian
pada pengembangan sistem \emph{plugin}, \emph{chunker}, \emph{indexer},
\emph{knowledge graph}, dan prapemrosesan kueri sebagai kontribusi utama.
Pembatasan ini juga membedakan komponen yang dievaluasi dari komponen pendukung
serta menetapkan konteks korpus dan eksperimen yang menjadi dasar penilaian
ketercapaian tujuan. Adapun batasan penelitian yang ditetapkan adalah sebagai
berikut.

\begin{enumerate}
    \item Omni-RAG dirancang sebagai arsitektur \emph{multi-domain}, dengan evaluasi utama menggunakan korpus dokumen hukum pendidikan di Indonesia dan dokumen \emph{customer service} Netgear.
    \item Kontribusi penelitian berfokus pada sistem \emph{plugin}, \emph{chunker}, \emph{indexer}, \emph{knowledge graph}, dan prapemrosesan kueri yang meliputi \emph{guardrail}, \emph{context PII}, serta \emph{history compactor}. Komponen lain dikerjakan oleh anggota penelitian yang berbeda.
    \item Sistem \emph{plugin} hanya ditujukan untuk memungkinkan pengguna melakukan kustomisasi logika komponen melalui \emph{adapter} tanpa mengubah kode inti. Kebutuhan yang berkaitan dengan perubahan infrastruktur, akses basis data, atau integrasi eksternal berada di luar cakupan \emph{plugin} dan dapat memerlukan perubahan pada kode inti.
    \item Representasi dokumen yang digunakan dalam penelitian hanya berbentuk teks. Gambar dan tabel tidak menjadi bagian dari representasi yang diproses.
    \item Aspek keamanan tidak mencakup perancangan keamanan infrastruktur, autentikasi, otorisasi, pengujian penetrasi, pemenuhan standar kepatuhan, atau aspek keamanan menyeluruh lainnya. Pembahasan dibatasi pada perlindungan PII dan pemeriksaan \emph{guardrail} dalam prapemrosesan RAG.
    \item \emph{PII}, \emph{guardrail}, dan pemadatan histori bukan fokus utama penelitian. Ketiganya hanya berfungsi sebagai komponen pendukung kebutuhan RAG secara umum dan tidak menjadi objek eksperimen.
\end{enumerate}

\section{Metodologi Penelitian}

Pelaksanaan penelitian ini menempuh beberapa tahapan umum sebagai berikut:

\begin{enumerate}
    \item Mengidentifikasi permasalahan dan kebutuhan sistem berdasarkan latar belakang, rumusan masalah, serta kondisi implementasi yang telah tersedia. Tahap ini menetapkan fokus kontribusi penelitian dan ruang lingkup komponen yang dikaji.
    \item Merancang arsitektur Omni-RAG beserta hubungan antarkomponennya, terutama \emph{plugin}, \emph{chunker}, \emph{indexer}, dan \emph{knowledge graph}. Rancangan juga mencakup komponen pendukung prapemrosesan kueri.
    \item Mengimplementasikan komponen yang telah dirancang dan melakukan verifikasi melalui pengujian serta pemeriksaan integrasi. Tahap ini memastikan setiap komponen dapat digunakan sesuai perannya dalam alur sistem.
    \item Menyiapkan data, konfigurasi, dan skenario eksperimen, kemudian menguji beberapa alternatif \emph{chunker}, \emph{indexer}, dan \emph{knowledge graph} pada kondisi yang ditetapkan.
    \item Menganalisis hasil eksperimen berdasarkan ukuran keberhasilan yang telah dirumuskan.
    \item Menarik kesimpulan berdasarkan hasil analisis untuk menentukan ketercapaian tujuan penelitian serta merumuskan peluang pengembangan selanjutnya.
\end{enumerate}

\section{Sistematika Pembahasan}

Laporan tugas akhir ini disusun ke dalam lima bab. Bab I memuat latar belakang,
rumusan masalah, tujuan beserta ukuran keberhasilannya, batasan penelitian, dan
metodologi yang menjadi dasar pelaksanaan penelitian. Bab II menguraikan kajian
pustaka yang menjadi landasan pemilihan arsitektur, sistem \emph{plugin}, unit
dokumen, strategi pengindeksan, \emph{knowledge graph}, pengelolaan konteks, dan
metode evaluasi. Bab III membahas analisis masalah serta rancangan solusi
Omni-RAG, termasuk arsitektur komponen, kontrak \emph{plugin}, \emph{lifecycle},
alur pemrosesan dokumen, alur kueri, dan rancangan penyimpanan. Bab IV menjelaskan
rancangan eksperimen atas komponen \emph{chunker}, \emph{indexer}, dan
\emph{knowledge graph}, menyajikan hasilnya, serta mengevaluasi ketercapaiannya
terhadap ukuran keberhasilan yang telah dirumuskan. Bab V menyimpulkan keterjawaban
rumusan masalah dan menyampaikan saran bagi pengembangan selanjutnya. Spesifikasi
kebutuhan, rancangan sistem, dan pengujian tambahan disajikan pada lampiran.
